Given an integer, your task is to find the number, sum and product of its divisors. As an example, let us consider the number $12$:

\textbullet\ the number of divisors is $6$ (they are $1$, $2$, $3$, $4$, $6$, $12$)
\textbullet\ the sum of divisors is $1+2+3+4+6+12=28$
\textbullet\ the product of divisors is $1 \cdot 2 \cdot 3 \cdot 4 \cdot 6 \cdot 12 = 1728$

Since the input number may be large, it is given as a prime factorization.

\vspace{0.3cm}\noindent\textbf{Input}

The first line has an integer $n$: the number of parts in the prime factorization.
After this, there are $n$ lines that describe the factorization. Each line has two numbers $x$ and $k$ where $x$ is a prime and $k$ is its power.

\vspace{0.3cm}\noindent\textbf{Output}

Print three integers modulo $10^9+7$: the number, sum and product of the divisors.

\vspace{0.3cm}\noindent\textbf{Constraints}

\textbullet\ $1 \le n \le 10^5$
\textbullet\ $2 \le x \le 10^6$
\textbullet\ each $x$ is a distinct prime
\textbullet\ $1 \le k \le 10^9$